\section{证据}
\begin{enumerate}
  \item 考得浅。
  \item 3 性。\zx{客观性/真实性、关联性、合法性}。\al{证人与当事人的利害关系 \yjt 客观性}。
  \item \zx{也会偶尔考直接、间接证据}。\al{当事人未能提供原件，不能单独作为定案依据 \dui}。
  \item 种类。会直接出题。
        \begin{enumerate}
          \item 书证。原则\zd{原件}。\zd{复制品/复印件 + 印章}（翻译机构\zd{或}个人，鉴定机构\zx{和}个人）。
          \item 物证。原则\zd{原件}。特殊证据可\zd{进行拍照、录像}。\al{毒品}。
          \item 视听资料。注明\zd{制作方法、制作人、制作日期、证明对象 + 文字记录说明 }。
          \item 电子数据。\al{聊天纪录}。
          \item 证人证言。\zd{签字} + \zx{身份证明}。\zd{费用败诉方承担}。\al{甲的证词有甲的签字后，即可作为证言使用 \cuo}。
          \item 勘验笔录。\zx{依申请 \hz 依职权}勘验现场。
                \begin{enumerate}
                  \item 出示法院证件。
                  \item 邀请基层组织 \hz 当事人所在单位派人参加。
                  \item 当事人 \hz 成年亲属应在场。不到场笔录说明情况即可。
                \end{enumerate}
          \item \zx{现场笔录}。
                \begin{enumerate}
                  \item 执法现场、制作主体是行政机关、\zx{签名或盖章}。\al{由经办人员签名并盖章 \cuo}。
                  \item \zx{拒绝签名 \yjt 注明即可}。不在场 \yjt 请见证人。二者\zx{互斥}。\al{必须有执法人员和甲的签名 \cuo}。
                \end{enumerate}
          \item 当事人陈述。
          \item 涉外证据。见三国。
          \item 考原告 vs 被告证据证明力问题。\zx{白纸黑字 > 说的}。
                \begin{enumerate}
                  \item \al{市卫生局提供的现场笔录的效力，优于某公司的证人对现场的描述 \dui}。
                  \item \al{执法局执法人员甲的证言优于乙（当地居民）的证言 \cuo}。
                \end{enumerate}
        \end{enumerate}
  \item \zx{举证责任}。\zx{不利被告原则}。
        \begin{enumerate}
          \item 原告。有举证权利而无需承担举证责任。例外。
                \begin{enumerate}
                  \item 符合起诉条件 \yjt 证明\zd{官管民 3 要素}。
                  \item 行政不作为存在 \yjt 证明\zd{申请过}。
                  \item 赔偿、补偿 \yjt 证明\zd{损害事实、结果}。
                  \item \bs{因被告的原因导致原告无法证明的 \yjt 被告}。要能反应过来啊！
                  \item \zd{超出合理范围} \yjt 原告。
                \end{enumerate}
          \item \bs{复议维持}。\bs{原行为} \yjt 原机关、复议机关\bs{共同}证明。\zd{复议决定} \yjt \zd{复议机关自己}承担。
        \end{enumerate}
  \item 证据补充。
        \begin{enumerate}
          \item 原告。原则\zd{可以补}充。例外。搞\zd{证据突袭，不予采纳}。
          \item 被告。原则\zd{禁止补}充。例外。第三人拿出\zd{新证}据。
        \end{enumerate}
  \item 证据调取。\yi \zx{依职权。国社三身程}。\er 依申请。\zd{被告不能申请}。\al{涉秘}。救济 \yjt \zd{复议}。
  \item 责令提交。\yi 依职权。\er 依申请 = \zd{开庭审理前 + 预付费用 + 书面}。没考过。
  \item 举证期。
        \begin{enumerate}
          \item 原告。原则\zd{开庭前}，最迟\zd{庭前证据交换前}。
          \item 被告。\zd{收到}起诉状副本 \zd{15} 日 + 延期情况消除后 \zd{15} 日。
        \end{enumerate}
  \item 质证。
        \begin{enumerate}
          \item 依职权调取的证据叫\zd{听取意见}（\al{质证 \cuo}）。\zd{涉密叫不公开质证}（\al{不质证 \cuo 任何证据都要进行质证}）。
          \item 证人要\zd{有行为能力}。可以要求到庭接受询问、签署保证书。
          \item \zd{有条件，可以}全程同步录音录像。
        \end{enumerate}
        \begin{kk}{证据的采信。证不真、路不正、来得晚}
          \begin{enumerate}
            \item \zd{没效力}。
            \item \zd{片面效力}。\zx{当时证据，当时主张}。
                  \begin{enumerate}
                    \item \al{甲诉讼中提供的、被告在复议中未作出处理的证据，可以作为认定被诉处理决定合法的依据 \cuo}。
                    \item 例外。\zd{复议机关为了认定原行为和复议决定合法而搜集和补充的}。
                  \end{enumerate}
            \item \zd{弱效力}。怀疑但无法否认 \yjt 证据补强规则。
          \end{enumerate}
        \end{kk}
  \item \zx{重新鉴定}。\zx{举证期限届满前}、自行 \hz \zx{书面申请}复议机关、费用均自担。
        \begin{enumerate}
          \item \al{当事人申请法院重新鉴定可以口头提出 \cuo}。
          \item \al{当事人申请法院重新鉴定应当再举证期限内提出 \dui}。
        \end{enumerate}
\end{enumerate}
